\section{Ampleness of every member of the family}\label{sec:ampleness}

For each $m$, we construct enough symmetric-power sections to embed
$\PP(E_m)$. The two quotient sequences are used alternately: their
kernels individually have one negative degree, but their tensor product
is positive in both rulings. The resulting filtration gives an explicit
very-ampleness range within Hartshorne's symmetric-power framework
\cite[Corollary~2.5 and Proposition~3.2]{Hartshorne1966}.

\begin{proposition}\label{prop:ampleness}
For every positive integer $m$, the bundle
$E_m=f_m^*G\otimes P$ is ample. More precisely,
$\OO_{\PP(E_m)}(n)$ is very ample for every even integer $n$ with
$n-1\geq3m$.
\end{proposition}

\begin{proof}
Fix $m$ and write
\[
 V=f_m^*G,\qquad \widetilde L_i=f_m^*L_i,\qquad
 \widetilde D=f_m^*D,\qquad
 K_i=\widetilde L_i(-\widetilde D),\qquad K=K_1\otimes K_2.
\]
Pullback of~\eqref{eq:quotient-sequences} gives
\begin{equation}\label{eq:pullback-quotients}
 0\longrightarrow K_i\longrightarrow V\longrightarrow
 \widetilde L_j\longrightarrow0,
 \qquad\{i,j\}=\{1,2\}.
\end{equation}
These sequences remain exact because their original quotients
are locally free. Here
\[
 K_1=\OO_S(21m,-3m),\qquad K_2=\OO_S(-3m,21m),
\]
and hence
\begin{equation}\label{eq:kernel-product}
 K=\OO_S(18m,18m),\qquad
 \OO_S(\widetilde D)=\OO_S(3m,3m).
\end{equation}
We write powers of line bundles for tensor powers throughout this
proof.

\smallskip\noindent\textbf{Alternating the quotient sequences.}
For an extension $0\to K'\to V'\to L'\to0$ with $K',L'$ line
bundles, multiplication gives
\[
 0\longrightarrow K'\otimes\Sym^{r-1}V'
   \longrightarrow\Sym^r V'\longrightarrow(L')^r
   \longrightarrow0\qquad(r\geq1).
\]
Indeed, a local splitting gives a kernel generator $k'$ and a lift
$e'$ of a quotient generator. The monomials containing $k'$ form the
kernel of the symmetric quotient, and the remaining monomial $(e')^r$
maps to a quotient generator. This is the usual symmetric-power
filtration of an extension; compare the proof of
\cite[Corollary~3.4]{Hartshorne1966}.

Apply this construction first to $V\to\widetilde L_2$ and then to
$V\to\widetilde L_1$. For $r\geq2$ it gives
\begin{align}
 0&\longrightarrow K_1\otimes\Sym^{r-1}V
   \longrightarrow\Sym^r V\longrightarrow\widetilde L_2^{\,r}
   \longrightarrow0,\label{eq:symmetric-first}\\
 0&\longrightarrow K\otimes\Sym^{r-2}V
   \longrightarrow K_1\otimes\Sym^{r-1}V
   \longrightarrow K_1\otimes\widetilde L_1^{\,r-1}
   \longrightarrow0.\label{eq:symmetric-second}
\end{align}
Thus each pair of steps leaves the residual bundle
$K\otimes\Sym^{r-2}V$, with one line-bundle quotient at each step.
For even $n$, iteration gives a filtration of $\Sym^nV$ with line
factors
\begin{equation}\label{eq:graded-lines}
 \begin{gathered}
 K^j\otimes\widetilde L_2^{\,n-2j},\qquad
 K^j\otimes\widetilde L_1^{\,n-2j}(-\widetilde D),
 \qquad 0\leq j<n/2,\\
 K^{n/2}.
 \end{gathered}
\end{equation}
Indeed, after $j$ pairs the residual bundle is
$K^j\otimes\Sym^{n-2j}V$, whose second quotient is
$K^jK_1\widetilde L_1^{\,n-2j-1}
 =K^j\widetilde L_1^{\,n-2j}(-\widetilde D)$.
Only the second quotient contains an additional twist by
$-\widetilde D$ beyond the factors absorbed into $K^j$.

Each step uses just one of the locally split sequences
\eqref{eq:pullback-quotients}, so it remains valid when the two kernel
lines have the same fiber image. Concretely, in the pulled-back basis
$e,f$ of~\eqref{eq:local-basis}, the kernels are generated by
$k=de-hf$ and $f$, with $h$ a unit. For $n=2$ the filtration is
\[
 \OO(kf)\subset kV\subset\Sym^2V.
\]
Since $e,k$ are also a basis of $V$, the module $kV$ has basis
$ke,k^2$ and is a direct summand of $\Sym^2V$, with complement
generated by $e^2$. Moreover,
\[
 kf=\frac{d}{h}ke-\frac1h k^2.
\]
The coefficient $-1/h$ is a unit, so $kf$ extends to a basis of $kV$.
Both inclusions therefore have locally free quotients.
This verification uses no smoothness of $d=0$ and applies at nodes
and to nonreduced pullbacks of the divisor.

\smallskip\noindent\textbf{Global generation after twisting.}
Choose $n$ even with $n-1\geq3m$, and set
\begin{equation}\label{eq:generated-symmetric-power}
 W=\Sym^n V\otimes P^{n-1}.
\end{equation}
By~\eqref{eq:kernel-product}, the bidegrees of its filtration
factors, in the order of~\eqref{eq:graded-lines}, are
\begin{align*}
 &(18mj+n-1,\;18mj+24m(n-2j)+n-1),\\
 &(18mj+24m(n-2j)-3m+n-1,\;18mj-3m+n-1),\\
 &(9mn+n-1,\;9mn+n-1).
\end{align*}
Every entry is nonnegative: the only possibly negative contribution
is $-3m$, and it is offset by $n-1$.

For $c,t\geq0$, products of homogeneous monomials generate
$\OO_S(c,t)$. Moreover $H^1(\PP^1,\OO(c))=H^1(\PP^1,\OO(t))=0$, as
computed on the two standard affine charts, so the K\"unneth formula
gives $H^1(S,\OO_S(c,t))=0$
\cite[Tags~01XD and~0BED]{StacksProject}.
If $0\to U\to Z\to Q\to0$ is an extension of globally generated
bundles with $H^1(S,U)=H^1(S,Q)=0$, the cohomology sequence gives
$H^1(S,Z)=0$ and lifts every global section of $Q$ to $Z$.
The lifted generators of $Q$, together with generators of $U$,
span each fiber of $Z$. Starting with the innermost line factor and
applying this observation successively proves that $W$ is globally
generated.

\smallskip\noindent\textbf{The projective embedding and its line bundle.}
We record the tautological twists in the quotient-projectivization
and Segre constructions
\cite[Tags~01OA, 01N1, and~01WD]{StacksProject}.
Let $\pi:\PP(V)\to S$ and $H=H^0(S,W)$. Evaluation
$H\otimes\OO_S\to W$ induces a closed immersion
\[
 \PP(W)\hookrightarrow\PP(H)\times S
\]
whose projective hyperplane line pulls back to
$\OO_{\PP(W)}(1)$. The relative Veronese closed immersion
$\PP(V)\to\PP(\Sym^n V)$ sends a quotient line $\ell$ to
$\ell^{\otimes n}$. Under the natural identification
$\PP(\Sym^n V)\simeq\PP(W)$, its tautological pullback is
\[
 \OO_{\PP(V)}(n)\otimes\pi^*P^{n-1}.
\]
Consequently their composition embeds $\PP(V)$ into
$\PP(H)\times S$ with this hyperplane pullback from the first
factor. Embedding $S$ by $P=\OO_S(1,1)$ and then applying the Segre
embedding gives a closed immersion into projective space whose
hyperplane line pulls back to
\[
 \bigl(\OO_{\PP(V)}(n)\otimes\pi^*P^{n-1}\bigr)
       \otimes\pi^*P
 =\bigl(\OO_{\PP(V)}(1)\otimes\pi^*P\bigr)^n.
\]
This line is therefore very ample. Finally,
$\PP(E_m)=\PP(V\otimes P)\simeq\PP(V)$ identifies
$\OO_{\PP(E_m)}(1)$ with
$\OO_{\PP(V)}(1)\otimes\pi^*P$, proving the proposition.
\end{proof}
